\section{Conclusion and Future Work}
\label{sec:conclusion}

\subsection{Conclusion}
\label{subsec:conclusion_summary}
We presented a hybrid robotic framework organized around thin-film materials experimentation, with heterogeneous manipulation and mobile transport coordinated by a failure-aware hierarchical state machine. Its method focus is simulation-based skill preparation and few-shot sim-to-real adaptation for precision liquid operations, combining behavioral cloning, domain randomization, and parameter-regularized fine-tuning. Posture-aware grasping and constraint-aware transport support workflow execution, while physical-to-simulation state synchronization serves as an auxiliary monitoring interface.

The reported individual-skill trials yield an 82.7\% pooled success rate for domain randomization with few-shot adaptation. This result concerns the evaluated skill trials, including the serial-dilution extension, and does not establish end-to-end thin-film reliability, volumetric precision, or performance under held-out conditions. Thin-film fabrication experiments have been completed, but the manuscript still requires quantitative material-quality comparisons before claiming improved uniformity or corrosion protection. The central application question remains whether gains in precision operation translate into more reliable workflows and higher accepted-film yield.

\subsection{Limitations and Future Work}
\label{subsec:limitations_future}
The priority is to complete the evidence linking operation performance to materials outcomes. This includes documenting run-level workflow records, first-pass and recovery-assisted success, completion time, and operator interventions; comparing manual, fixed-trajectory, and learned execution using independent substrates; and reporting thickness uniformity, contact angle, roughness, and corrosion measurements with uncertainty estimates. Direct spill measurements are required to support liquid-retention claims beyond trajectory smoothness.

For transfer evaluation, the supplementary protocol compares training variants and real-demonstration budgets using independent test runs, precision measurements, and held-out conditions. Multiple training seeds and adaptation subsets are needed to assess variability and avoid attributing batch-specific differences to the method. Demonstration acquisition remains a manual cost. State-synchronization summaries also require raw trajectories and independent references for verification, but synchronization is not the primary contribution. Extension to serial dilution is secondary to validating the thin-film workflow and does not establish assay quality or contamination control without corresponding measurements.
